\section{The Fused Multi-Shell Decoder}
\label{sec:decoder}

\subsection{Background: what a $\Lambda_{24}$ index is}

LLVQ quantizes weights in blocks of 24: each is normalized and its direction
snapped to the nearest point of the Leech lattice $\Lambda_{24}$, within the
ball $\Lambda_{24}(12)$ (squared norm $\le 2\cdot 12$ in the standard
scaling). The ball holds $N = 111\,043\,117\,458\,000$ points
($70\,486\,236\,999\,360$ in its outermost shell), so a direction index
costs $\lceil\log_2 N\rceil = 47$ bits; one gain bit selects between two
per-tensor magnitude centroids. Total: 48 bits per 24-weight block---exactly
2 bits/weight of payload.

The 47 bits are close to tight: empirical index entropy over the shipped
model's $150$M blocks is $46.6536$ bits against 47 paid, a $0.74\%$ margin
that entropy-coding the rank would close.

In $\sqrt{8}\,\Lambda_{24}\subset\mathbb{Z}^{24}$, lattice points are
characterized by three simultaneous conditions tied to the binary Golay code
$\mathcal{G}_{24}$ \citep{conway1999sphere}: parity (coordinates all even or
all odd), Golay (coordinates in a fixed residue class mod~4 support a
codeword), and a sum condition mod~8. Every lattice point \emph{in the ball}
falls in one of 301
\emph{classes}---orbits under coordinate permutation and admissible sign
changes, each defined by a Golay codeword weight, a magnitude multiset, and a
coset (even/odd). Format~v1 enumerates classes, then arrangements and signs
within a class, in fixed order; the 48-bit-per-block index is bijective and
pinned by tests down to its mixed-radix composition orders.

\subsection{The encoder/decoder asymmetry}

One asymmetry governs everything: the \emph{encoder} (nearest-neighbor search
over $1.1\times10^{14}$ points) runs offline once per model and may cost
hours; the \emph{decoder} (index~$\to$~vector) runs in every matvec and must
cost shifts and masks. The original paper's kernel decodes a single shell
($m=3$); at the 2-bit rate the codebook spans eleven shells and 301 classes,
where a naive per-class branch would diverge across every warp.

Our kernel never decodes the 47-bit index at inference: an offline
\emph{transcoder} expands each block index once, at load time, into a
decoded-field representation the kernel reads directly
(\S\ref{sec:layouts}). \textsc{Slot32}, the simplest form, holds a 9-bit
class identifier, the gain bit, a 24-bit sign mask, and four 24-bit
coordinate masks selecting the coordinates carrying each magnitude level.
Reconstructing coordinate $i$ is a table lookup of the class's levels, a mask
test per level, and a sign flip---the same instruction sequence for every
lane whatever the class, hence zero warp divergence.
The per-class table is small---512 records of 24 bytes (five level values and
a level count), 12~KB---and sized at 512 rather than the class count so a
9-bit class field cannot index out of bounds. It is read through L1 as an
ordinary \texttt{const}-qualified global pointer, not CUDA constant memory, which by construction serializes the 32 distinct addresses a warp
issues here; \S\ref{sec:attribution} bears the choice out (0.041~ms, the
smallest term measured). Every point of a shell has the same norm, so the
normalization between lattice integer coordinates and unit direction is a
per-class constant folded into the stored level values: the union-of-shells
rescaling that the original paper's Appendix~G flags as the hardware cost of
a multi-shell codebook reduces to one multiply per block.

\subsection{Verification discipline}

Lattice decoders fail silently: a wrong sign convention or swapped
enumeration order yields plausible weights and a ruined model. Three defenses:

\begin{itemize}
\item \textbf{Bit-exact transcoding.} The transcoder expanding format~v1
into the four distinct byte formats of \S\ref{sec:layouts} (the two
\textsc{Golay70} rows share one, differing only in how the kernel decodes it)
is verified bit-exact across all layouts, its test suite hardened by
mutation testing (${\sim}25$ mutants killed). A surviving mutant means a
weak test or dead code; both stop the line.

\item \textbf{Line-by-line numerical verification before any timing.} Every
benchmark first checks every output row of every arm against an independent
f64 reference ($1\,105\,920$ rows for the full-model benchmark), threshold
$10^{-5}\cdot\sum|w\cdot x|$; our six arms come in at
$2.2$--$3.4\times10^{-8}$, consistent with f32 accumulation. The 4-bit
competitor stores a binary16 output, so is held to a matching $10^{-3}$
threshold ($5.8\times10^{-5}$ measured) against a reference still exact bit
for bit. No arm is timed until proven to compute the right thing.

\item \textbf{Invariant locks upstream.} The lattice code is pinned by exact
combinatorial invariants: the kissing number $196\,560$, the theta-series
coefficients, and the class-cardinality sum
$N(13)=280\,974\,212\,784\,720$, which no incorrect enumeration constraint
can reproduce.
\end{itemize}
